\chapter{Software implementation}
\label{ch:implementation}

This chapter presents the software implementation of the unified benchmarking apparatus described in the preceding chapter, together with the experimental results it produces.  Section~\ref{sec:impl-modules} walks through the key source modules that realise the design of Chapter~\ref{ch:method}.  Section~\ref{ch:results-detection} then reports the same-dataset detection benchmark and metric-sensitivity analysis on HAI~21.03 and Morris.  Section~\ref{ch:results-transfer} reports the cross-testbed transfer study under both source-calibrated and target-calibrated thresholds and the within-HAI leave-one-process-out control.  Section~\ref{ch:results-attribution} reports the per-sensor attribution evaluation with multi-seed sensitivity, including SHAP attribution for the two classical baselines.  Section~\ref{ch:discussion} discusses the implications of the findings, anticipated objections, and limitations.

\section{Key implementation modules}
\label{sec:impl-modules}

The codebase is organised so that the design rules of Chapter~\ref{ch:method} are enforced in code rather than left to caller discipline.  This section presents the four modules that carry the most methodological weight: the detector interface that makes the six models interchangeable, the scaler-leak guard, the schema-alignment projection, and the point-adjust expansion.  The listings are abridged from the committed sources under \file{src/}.

\subsection{The unified detector interface}

Every model implements the abstract base class in Listing~\ref{lst:base}.  Three methods --- \code{fit}, \code{score}, \code{attribute} --- are the entire contract the experiment drivers rely on, which is what lets the same code path serve same-dataset detection, transfer, LOPO, and attribution.  \code{attribute} defaults to raising, so a model that cannot localize fails explicitly rather than silently returning a meaningless vector; \code{supports\_attribution} lets the attribution driver skip such models.

\begin{lstlisting}[language=Python, caption={The \code{AnomalyDetector} abstract base class (\file{src/models/base.py}).}, label={lst:base}]
class AnomalyDetector(ABC):
    name: str = "AnomalyDetector"

    @abstractmethod
    def fit(self, X_train, X_val=None) -> "AnomalyDetector":
        """Fit on normal-only data. Must never see attack rows."""
        ...

    @abstractmethod
    def score(self, X) -> np.ndarray:
        """Per-sample anomaly score; higher = more anomalous."""
        ...

    def attribute(self, X) -> np.ndarray:
        """Per-feature contribution; AE family and TranAD override this."""
        raise NotImplementedError(
            f"{type(self).__name__} does not implement attribute().")

    def supports_attribution(self) -> bool:
        return type(self).attribute is not AnomalyDetector.attribute
\end{lstlisting}

\subsection{The scaler-leak guard}

Listing~\ref{lst:scale} implements Equation~\eqref{eq:minmax}.  The first line of \code{scale\_bundle} asserts that no attack row reached the train or validation split; the scaler is then fit on the training split alone and used to transform all three splits.  Fitting before the assertion would already have leaked attack-range statistics, so the order is load-bearing.

\begin{lstlisting}[language=Python, caption={Leak-guarded min--max scaling (\file{src/preprocessing.py}).}, label={lst:scale}]
def scale_bundle(bundle: DatasetBundle) -> ScaledArrays:
    bundle.assert_no_attack_in_train_val()   # fail loudly before fitting
    scaler = MinMaxScaler(clip=True)
    scaler.fit(bundle.X("train"))            # statistics from normal data only
    return ScaledArrays(
        X_train=scaler.transform(bundle.X("train")).astype(np.float32),
        X_val=scaler.transform(bundle.X("val")).astype(np.float32),
        X_test=scaler.transform(bundle.X("test")).astype(np.float32),
        y_test=bundle.y("test"),
        feature_names=list(bundle.features.columns),
        scaler=scaler,
    )
\end{lstlisting}

\subsection{The schema-alignment projection}

Listing~\ref{lst:project} implements the type-aggregation operator $\Pi_d$ of Equation~\eqref{eq:projection}.  Each canonical type becomes one output column, computed by aggregating all same-typed raw features of the dataset; a type present in the target schema but absent from this dataset yields a zero column, so the model always receives a fixed-width input over the type intersection.

\begin{lstlisting}[language=Python, caption={Type-vector projection for cross-testbed transfer (\file{src/transfer/schema\_align.py}).}, label={lst:project}]
def project_dataframe(df, feat_to_type, target_types, aggregations):
    type_to_feats = _type_to_features(feat_to_type)   # drops excluded types
    cols = {}
    for t in target_types:
        feats = [f for f in type_to_feats.get(t, []) if f in df.columns]
        if not feats:
            cols[t] = np.zeros(len(df), dtype=np.float32)   # missing type
            continue
        agg = aggregations.get(t, "mean")                   # max for binary
        block = df[feats].to_numpy(dtype=np.float32)
        cols[t] = _AGG_FNS[agg](block, axis=1)
    return pd.DataFrame(cols, columns=target_types).astype(np.float32)
\end{lstlisting}

\subsection{The point-adjust expansion}

Listing~\ref{lst:pa} implements Equation~\eqref{eq:pa-expand}.  It walks the ground-truth series, and for every contiguous attack event in which at least one timestep was predicted positive, it flips the whole event to positive before the point-wise F1 is computed.  Isolating the expansion in one auditable function is what lets the unit tests pin the inflation effect that Section~\ref{ch:results-detection} reports.

\begin{lstlisting}[language=Python, caption={Point-adjust prediction expansion (\file{src/evaluation/point\_adjust.py}).}, label={lst:pa}]
def _expand_predictions(y_true, y_pred):
    out = y_pred.copy()
    in_attack, start = False, 0
    for t in range(len(y_true)):
        if y_true[t] == 1 and not in_attack:
            in_attack, start = True, t
        elif y_true[t] == 0 and in_attack:
            if y_pred[start:t].any():        # any hit -> credit whole event
                out[start:t] = 1
            in_attack = False
    if in_attack and y_pred[start:].any():
        out[start:] = 1
    return out
\end{lstlisting}

Together these four modules fix the methodological invariants in code: models are interchangeable behind one interface, scaling cannot leak, the cross-testbed projection has a single definition, and the most inflation-prone metric is isolated for testing.  The remaining sections of this chapter report what the apparatus produces.

\chapter{Results I --- Same-dataset detection and metric sensitivity}
\label{ch:results-detection}

Notebooks: \file{04\_sota\_results.ipynb}, \file{06\_metric\_sensitivity.ipynb}.  Source of numbers: \file{results/metrics/summary.parquet}, same-dataset rows.

\section{Full detection grid}

\begin{table}[h]
\centering
\caption{Detection F1 across six models, two datasets, three metrics.  Best per (dataset, metric) in bold.}
\label{tab:det-grid}
\small
\begin{tabular}{lrrrrrr}
\toprule
 & \multicolumn{3}{c}{HAI} & \multicolumn{3}{c}{Morris} \\
\cmidrule(lr){2-4}\cmidrule(lr){5-7}
Model & eTaPR & PA-F1 & pointwise & eTaPR & PA-F1 & pointwise \\
\midrule
Isolation Forest & 0.243 & \textbf{0.621} & 0.185 & 0.145 & 0.297 & 0.090 \\
One-Class SVM    & 0.170 & 0.540 & 0.265 & \textbf{0.369} & 0.631 & \textbf{0.275} \\
Dense AE         & 0.057 & 0.477 & \textbf{0.356} & 0.356 & 0.586 & 0.253 \\
LSTM-AE          & 0.147 & 0.221 & 0.105 & --- & --- & --- \\
USAD             & 0.151 & 0.222 & 0.107 & 0.009 & 0.030 & 0.001 \\
TranAD           & \textbf{0.276} & 0.458 & 0.189 & 0.339 & \textbf{0.647} & 0.129 \\
\bottomrule
\end{tabular}
\end{table}

Figure~\ref{fig:det-grid} visualises the same data.

\begin{figure}[h]
\centering
\includegraphics[width=0.95\textwidth]{04_sota/sota_six_model_comparison.png}
\caption{Six-model detection F1 on HAI (left) and Morris (right) under three metrics.  Rankings change across metrics and across datasets.}
\label{fig:det-grid}
\end{figure}

\section{Ranking changes across metrics}
\label{sec:rank-change}

The rankings are not stable:
\begin{itemize}
  \item On HAI: best eTaPR $=$ TranAD; best PA-F1 $=$ Isolation Forest ($0.621$); best pointwise $=$ Dense AE.
  \item On Morris: best eTaPR $=$ OC-SVM ($0.369$); best PA-F1 $=$ TranAD ($0.647$); best pointwise $=$ OC-SVM.
\end{itemize}

Using only one metric, a reader could rank any of OC-SVM, TranAD, Isolation Forest, or Dense AE as the ``winner'' and publish a confident title.  \textbf{This is why we report three.}

\section{PA-F1 inflation on classical detectors}

\begin{itemize}
  \item Dense AE HAI: pointwise $0.356$, PA-F1 $0.477$ ($1.34\times$ inflation).
  \item Isolation Forest HAI: pointwise $0.185$, PA-F1 $0.621$ ($3.35\times$ inflation).
  \item OC-SVM Morris: pointwise $0.275$, PA-F1 $0.631$ ($2.29\times$ inflation).
\end{itemize}

Classical models produce isolated high-score points rather than contiguous positive runs; PA-F1 grants full segment credit to those isolated flags, inflating the headline numbers disproportionately.  eTaPR penalises segment coverage and reduces the inflation.  This is the Kim~2022 \cite{kim2022towards} critique replicated on our data.

\section{USAD collapse on Morris}

USAD eTaPR on Morris $= 0.009$.  The model's test-time scores are effectively constant --- the discrete Modbus features break the USAD assumption of a continuous normal manifold.  Audibert et al.~\cite{audibert2020usad} did not evaluate on a Modbus RTU dataset; we flag this as a reviewer-facing pitfall for researchers extending USAD to protocol-level ICS data.

\section{HAI LSTM-AE reproducibility check}

Seed 42, val-percentile threshold 99, 10 epochs on a 30k window subsample: pointwise $0.105$, PA $0.221$, eTaPR $0.147$.  Shin et al.~\cite{shin2020hai} report $\sim 0.37$--$0.45$ eTaPR for LSTM baselines on HAI 20.07.  The gap is $\sim 0.2$ eTaPR; outside the $\pm 0.05$ project tolerance.  Sources of the gap: dataset version 20.07 vs.\ 21.03 (different attack set); 30k subsample of $\sim 1.5\mathrm{M}$ windows for CPU tractability; conservative threshold (99th percentile vs.\ oracle best-F1).  Per project policy, the gap is documented rather than tuned.  The qualitative ordering (windowed deep models $>$ classical on HAI eTaPR) is preserved.

\section{HAI TranAD reproducibility check}
\label{sec:tranad-repro}

Seed 42, RTX 3050 Ti GPU, 30 epochs, 30k window subsample: pointwise $0.189$, PA $0.458$, eTaPR $0.276$.  Tuli et al.~\cite{tuli2022tranad} report PA-F1 $\sim 0.97$ on HAI 20.07.  The gap is $\sim 0.5$ PA-F1 --- large.  Additional gap sources beyond \S4.5: single-layer encoder with $d_{\text{model}} = 64$ (our config) vs.\ the paper's deeper network; oracle-PA-best threshold in the paper.

The qualitative ranking still holds: TranAD is \#1 on HAI under PA-F1 in our setup.  The value of the result is in the comparison across models under a consistent protocol, not in matching published absolutes.

\chapter{Results II --- Cross-testbed transfer and within-HAI LOPO}
\label{ch:results-transfer}

This is the headline chapter.  Notebook: \file{05\_cross\_dataset.ipynb}.  Figures: F2, F3, F4 under \file{results/figures/05\_cross\_dataset/}.

\section{Schema alignment}

Six canonical types intersect HAI and Morris after hand-tagging in \file{data/feature\_types.yaml}: pressure, pump\_state, setpoint, system\_state, valve\_position, control\_signal.  HAI contributes 79 features collapsing to 6; Morris 16 features collapse to 6.  The projection is lossy by design --- we trade per-variable fidelity for a cross-testbed input shape that every detector can consume.

\section{Finding 3.1 --- source calibration collapses to the class-prior baseline}
\label{sec:finding-src-cal}

Under \code{src-cal} (threshold $= $ 99th percentile of source validation scores), every transferred detector degenerates to a predict-all-attack baseline calibrated to the \emph{target's} class prior.

\begin{table}[h]
\centering
\caption{Source-calibrated cross-testbed F1 tracks the target class prior rather than any detector property.}
\small
\begin{tabular}{lrrl}
\toprule
Target & Class prior (windowed) & Predict-all F1 & Observed F1 range \\
\midrule
Morris & 0.91 & 0.95 & 0.95--0.99 (windowed); 0.62--0.75 (classical) \\
HAI    & 0.03 & 0.05 & 0.04--0.07 (all models) \\
\bottomrule
\end{tabular}
\end{table}

The score distributions on source and target data have no calibrated overlap; a fixed percentile threshold either flags everything or nothing.  \textbf{A paper reporting only \code{src-cal} numbers on Morris (say F1$=0.98$) is reporting a class-prior artifact, not a transfer success.}

\begin{figure}[h]
\centering
\includegraphics[width=0.95\textwidth]{05_cross_dataset/transfer_same_vs_cross.png}
\caption{Same-domain target F1 vs.\ source-calibrated transfer F1 per model.  Transfer numbers mirror the target's predict-all baseline (grey); the observed F1 is not detecting anything.}
\label{fig:transfer-same-vs-cross}
\end{figure}

\section{Finding 3.2 --- target calibration reveals a classical-vs-SOTA split}
\label{sec:finding-tgt-cal}

Under \code{tgt-cal} (threshold from 99th percentile of \emph{target}-normal validation --- unlabeled, attack-free data available in any realistic deployment), the collapse disappears and a genuine separation emerges.

\begin{table}[h]
\centering
\caption{HAI $\rightarrow$ Morris transfer F1 under \code{tgt-cal}.  Windowed SOTA models separate from classical ones only once thresholds are target-calibrated.}
\small
\begin{tabular}{lrr}
\toprule
Model & tgt-cal eTaPR & tgt-cal PA-F1 \\
\midrule
Dense AE          & 0.07 & 0.14 \\
OC-SVM            & 0.08 & 0.17 \\
Isolation Forest  & 0.13 & 0.22 \\
\textbf{LSTM-AE}  & \textbf{0.44} & \textbf{0.81} \\
\textbf{TranAD}   & \textbf{0.38} & \textbf{0.77} \\
\textbf{USAD}     & \textbf{0.40} & \textbf{0.74} \\
\bottomrule
\end{tabular}
\end{table}

The classical three are clustered near zero; the windowed SOTA three sit at $0.38$--$0.44$ eTaPR.  This is the first crisp classical-vs-SOTA gap in our study and it is invisible under source calibration.

The reverse direction (Morris $\rightarrow$ HAI) is harder: pointwise F1 stays under $0.07$ for all models because the HAI $2.4\%$ attack prior saturates any 99th-percentile operating point.  PA-F1 recovers for classical detectors (Dense AE $0.47$, OC-SVM $0.36$) via the segment-inflation mechanism critiqued in Chapter~\ref{ch:results-detection} --- the same weakness that boosts classical PA-F1 same-dataset.

\begin{figure}[h]
\centering
\includegraphics[width=0.98\textwidth]{05_cross_dataset/transfer_calibration_comparison.png}
\caption{Source-cal vs.\ target-cal transfer F1 vs.\ same-domain baseline for each model.  Target calibration alone converts the trivial-baseline collapse into a meaningful measurement.}
\label{fig:transfer-cal-comp}
\end{figure}

\subsection{Why only the windowed models transfer}

Hypothesis: the classical models' normal manifold is essentially per-dataset specific; they fit the density of the source type-vector without capturing the temporal structure that is approximately shared between testbeds.  Windowed models encode short-term dynamics of normal operation --- pressure rising in ramp phases, valves opening after setpoints move --- and those dynamics recur in both testbeds.  The representation retains the invariant; the operating point does not.

\section{Finding 3.3 --- within-HAI LOPO isolates feature shift}
\label{sec:finding-lopo}

Holding class prior fixed (both train and test are HAI), we drop one process's feature columns and compute F1 per-attack-target-process.

\subsection{P2 drop produces targeted blindness}

\begin{table}[h]
\centering
\caption{F1 on P1-attacks vs.\ P2-attacks when P2 features are held out.  The drop is targeted to P2, not global.}
\small
\begin{tabular}{lrr}
\toprule
Model & F1 on P1-attacks (drop-P2) & F1 on P2-attacks (drop-P2) \\
\midrule
Dense AE & 0.41 & 0.03 \\
LSTM-AE  & 0.09 & 0.00 \\
USAD     & 0.28 & 0.07 \\
\bottomrule
\end{tabular}
\end{table}

P2 is the turbine controller; its sensors are almost all discrete state flags with no analogue in P1, P3, or P4.  Removing them leaves the model blind to attacks that manifest only through those flags --- while detection on P1 attacks barely moves.

\subsection{P1 drop produces global degradation}

When P1 (water loop, $33$ of $79$ features, $82\%$ of test attacks) is removed, all three models drop to F1 $\sim 0.1$ on every attack class.  P1 carries most of the model's picture of ``normal''; losing those features starves the whole detector.

\subsection{P3 and P4 drop is mild and sometimes helpful}

Dropping P3 or P4 occasionally \emph{improves} detection: USAD drop-P3 F1 on P1-attacks $= 0.57$ (vs.\ $0.35$ in the Phase-2 same-domain baseline).  Narrower input regularizes the windowed models.

\begin{figure}[h]
\centering
\includegraphics[width=0.95\textwidth]{05_cross_dataset/lopo_heatmap.png}
\caption{LOPO per-attack-process pointwise F1.  Rows are the held-out process; columns are the process the attack targets.  P4 column is empty because HAI 21.03 has no P4 attacks.}
\label{fig:lopo-heatmap}
\end{figure}

\subsection{No consistent SOTA advantage under LOPO}

Unlike the cross-testbed setting (where SOTA clearly beat classical under \code{tgt-cal}), LOPO shows no consistent classical/SOTA gap.  Dense AE is competitive with USAD on the P2-blindness diagnostic; both beat LSTM-AE on absolute numbers under drop-P2.  The cross-testbed SOTA advantage therefore looks tied to \emph{cross-schema representation}, not to robustness to feature loss per se.

\section{Cross-reading: sensor-level transfer ceilings}

Best within-HAI LOPO pointwise F1 $= 0.57$ (USAD, drop-P3, P1-attacks).  Best cross-testbed \code{tgt-cal} pointwise F1 $\approx 0.13$ (LSTM-AE HAI$\rightarrow$Morris).  Same architecture, same class-prior control, partial sensor overlap: the cross-dataset ceiling is an order of magnitude lower than the within-dataset one.  This defends the cross-testbed numbers against the ``you just need a better common feature space'' reviewer critique --- even under the best within-testbed conditions, sensor-level transfer loses half its detection capability.

\section{External-validation pipeline}
\label{sec:external-validation}

The transfer study above is closed over the two datasets chosen in Chapter~\ref{ch:method}.  To allow any later Morris-family capture (re-collected gas-pipeline ARFFs, the water-storage companion, future variants) to be scored against the same trained detectors without retraining, we ship a small external-validation pipeline: each run in \file{experiments/configs/} can opt into saving its fitted model, scaler, and threshold as a self-contained artifact, and a thin CLI (\code{python -m scripts.score\_external}) plus FastAPI app (\file{app/main.py}) apply that artifact to any uploaded ARFF/CSV.  The artifact reproduces the original summary-parquet row to floating-point tolerance on the held-out test split, so the pipeline is suitable both for the reproducibility gate (verifying no drift between train-time metrics and post-hoc scoring) and for extending the cross-testbed analysis beyond what is in this thesis.  The tooling is strictly Phase~A: gas-pipeline adapter only, localhost only, labels required for the metric panel; water-storage and power-system adapters are future work (\S\ref{sec:future}).

\section{Per-sensor attribution}
\label{ch:results-attribution}

Notebook: \file{07\_attribution.ipynb}.  Figure: F5.  Source: \file{results/metrics/attribution.parquet} ($3$ seeds $\times$ $4$ AE-family models $\times$ reconstruction; $3$ seeds $\times$ TranAD $\times$ attention) and \file{results/metrics/attribution\_shap.parquet} ($3$ seeds $\times$ $\{$IF, OCSVM$\}$ $\times$ SHAP) for the two classical baselines that don't have reconstruction-error attribution.

\subsection{Evaluation protocol recap}

For each attack window in HAI test, the model produces a per-feature attribution vector.  We rank features by attribution and compute precision@k --- the fraction of the top-k features that belong to the attacked process (from HAI's \code{attack\_P\{n\}} label).  Random baseline $=$ share of features in the attacked process: $0.48$ for P1, $0.28$ for P2, $0.09$ for P3.

\subsection{Precision@5 (mean $\pm$ std across seeds 7, 42, 123)}

\begin{table}[h]
\centering
\caption{Precision@5 by model and attacked process, averaged over three seeds.  Best per process in bold.}
\label{tab:attr-p5}
\small
\begin{tabular}{lrrr}
\toprule
Model & P1 p@5 & P2 p@5 & P3 p@5 \\
\midrule
Dense AE & $0.69 \pm 0.01$ & $0.74 \pm 0.00$ & $\mathbf{0.33 \pm 0.02}$ \\
LSTM-AE  & $0.74 \pm 0.00$ & $0.09 \pm 0.00$ & $0.20 \pm 0.00$ \\
USAD     & $0.74 \pm 0.00$ & $0.09 \pm 0.00$ & $0.21 \pm 0.02$ \\
TranAD   & $0.58 \pm 0.17$ & $\mathbf{0.35 \pm 0.19}$ & $0.02 \pm 0.01$ \\
Random   & $0.48$ & $0.28$ & $0.09$ \\
\bottomrule
\end{tabular}
\end{table}

\begin{figure}[h]
\centering
\includegraphics[width=0.95\textwidth]{07_attribution/attribution_p_at_5.png}
\caption{Attribution precision@5 by model and attacked process, with $\pm 1\sigma$ error bars across three seeds.  Dashed line is the random baseline for that process.}
\label{fig:attr-p5}
\end{figure}

\subsection{Finding 4.1 --- detection-vs-localization tradeoff (robust across seeds)}

Lift over random baseline for mean p@5:

\begin{table}[h]
\centering
\small
\begin{tabular}{lrrr}
\toprule
Model & P1 lift & P2 lift & P3 lift \\
\midrule
\textbf{Dense AE} & $\mathbf{1.45\times}$ & $\mathbf{2.68\times}$ & $\mathbf{3.72\times}$ \\
LSTM-AE & $1.55\times$ & $0.32\times$ & $2.22\times$ \\
USAD    & $1.55\times$ & $0.33\times$ & $2.33\times$ \\
TranAD  & $1.22\times$ & $1.26\times$ & $0.23\times$ \\
\bottomrule
\end{tabular}
\end{table}

\textbf{Dense AE --- the weakest Phase-2 detector --- is the only AE-family / windowed model whose attribution is reliably above random across all three processes.}  Seed variance is small ($\sim 0.01$--$0.02$ on P3); robust result, not seed artifact.  Finding~4.4 below shows the two classical baselines (IF and OCSVM via SHAP) join Dense AE in this everywhere-above-random group, with OCSVM in particular reaching Dense-AE-level localization despite its bottom-of-the-table detection F1.

The mechanism: Dense AE operates per-timestep, so per-feature reconstruction error is aligned with the current frame's anomaly.  Windowed models reconstruct a $60$-step window and average per-feature error over time; the averaging concentrates mass on features with the largest dynamic-range variation --- HAI's P1 flow and pressure sensors.  A small perturbation to a few discrete P2 flags is lost in the mean.

\textbf{This inverts the detection ranking of Section~\ref{ch:results-detection}}: TranAD is \#1 on HAI eTaPR but near-worst on P3 localization ($0.02$ mean p@5, $0.23\times$ random lift).  For an operator triaging a flagged window, Dense AE's per-feature output is more actionable.  Reporting only detection F1 obscures this.

\subsection{Finding 4.2 --- attention rollout is null (robust across seeds)}
\label{sec:finding-attn-null}

TranAD attention-weighted attribution differs from plain reconstruction attribution by $< 0.01$ on every $(\text{process}, k)$ pair across all three seeds.  The two methods are empirically interchangeable.

Intuition: with a single-layer encoder, rollout $=$ raw attention matrix.  Attention averaged over $60$ output positions on HAI normal data is close to uniform; reweighting timesteps uniformly leaves the per-feature ranking unchanged.  Attention-based XAI for per-sensor localization in ICS windowed AD is not automatically the right tool --- one has to verify attention-over-time is structurally variable before investing.

\subsection{Finding 4.3 --- LSTM-AE $\equiv$ USAD attribution is a genuine structural result}

Across $3$ seeds, LSTM-AE and USAD produce p@k values identical to two decimals on every $(\text{process}, k)$ pair.  Standard deviations are $\leq 0.016$ within each model.  This is not seed noise; the two architectures converge to the same per-feature reconstruction ranking on HAI.

Hypothesis: the ranking is dominated by the MinMax-scaled feature variance structure of HAI's $79$ sensors.  P1 sensors carry the highest dynamic range and feature count, so any architecture with enough capacity to minimise windowed reconstruction error learns to prioritise them.  LSTM-AE's sequence objective and USAD's adversarial dual-head objective both converge to the same optimum in the feature-ranking sense.

If your model class produces the same attribution ranking as a simpler baseline, the extra architectural complexity is buying detection accuracy (USAD and LSTM-AE \emph{do} differ in detection F1 slightly), not attribution quality.  Operators paying the complexity tax for a richer model should audit that they're getting the richer output.

\subsection{Revised framing vs.\ the seed-42-only report}

An earlier draft (commit \texttt{a2957b0}) reported ``TranAD actively anti-localizes on P3 (p@5 $= 0.02$, lift $0.2\times$, worse than flipping coins).''  Multi-seed averaging softens this claim:
\begin{itemize}
  \item P3 remains a TranAD weakness (mean $0.02$, range $[0.02, 0.03]$ across seeds --- consistent, not seed-random).
  \item But TranAD on P2 had large seed variance: p@5 $= 0.23$ at seed $42$ vs.\ $0.57$ at seed $7$.  Mean $0.35$, \emph{above} the $0.28$ random baseline.
\end{itemize}

Revised framing: TranAD \emph{can} localize to P2 with enough seed-trial budget, but its localization is seed-unstable, whereas Dense AE and LSTM-AE/USAD are seed-stable (even if the latter two are structurally stuck on P1).  An argument for multi-seed evaluation in any attribution study.

\subsection{Finding 4.4 --- SHAP fills in the classical baselines}
\label{sec:finding-shap}

The four AE-family / windowed models above use reconstruction-error or attention-rollout attribution, both of which require an internal reconstruction signal.  Isolation Forest and One-Class SVM produce a single anomaly score per row with no per-feature decomposition, so the same evaluation cannot be applied directly.  We close that gap with SHAP: \code{shap.TreeExplainer} on the underlying \code{IsolationForest} ensemble (exact, $\sim$ms/row), and \code{shap.KernelExplainer} on the \code{OneClassSVM} decision function with a deterministic $100$-row background sample drawn from the training-normal split (approximate, $\sim$seconds/row).  The KernelExplainer's predictor wraps \code{-decision\_function} so SHAP explains the project's anomaly-score convention (higher contribution = more anomalous) directly, and the background sample is fixed across the three training seeds so SHAP variance reflects the artifacts, not the sampling.  Evaluation reuses the same \code{precision\_at\_k\_by\_attack} harness from \S\ref{ch:results-attribution}; the only deviation from Phase 4's protocol is that OCSVM evaluation subsamples to $100$ attack rows per process per seed to keep KernelExplainer's runtime bounded --- the precision@k aggregate is robust to this at $n_{\text{attack}} \geq 100$.

\begin{table}[h]
\centering
\caption{SHAP-based precision@$5$ for the two classical baselines, mean over seeds $\{7, 42, 123\}$, alongside the AE-family numbers from Table~\ref{tab:attr-p5}.  Random baselines: $0.48$ for P1, $0.28$ for P2, $0.09$ for P3.}
\label{tab:attr-shap-p5}
\small
\begin{tabular}{llrrr}
\toprule
Model & Method & P1 p@5 & P2 p@5 & P3 p@5 \\
\midrule
Isolation Forest & SHAP-tree   & $0.67 \pm 0.03$ & $0.38 \pm 0.04$ & $0.19 \pm 0.06$ \\
OCSVM            & SHAP-kernel & $0.69 \pm 0.02$ & $\mathbf{0.42 \pm 0.03}$ & $0.25 \pm 0.00$ \\
\midrule
Dense AE & reconstruction & $0.69$ & $0.74$ & $\mathbf{0.33}$ \\
LSTM-AE  & reconstruction & $0.74$ & $0.09$ & $0.20$ \\
USAD     & reconstruction & $0.74$ & $0.09$ & $0.21$ \\
TranAD   & reconstruction & $0.58$ & $0.35$ & $0.02$ \\
\bottomrule
\end{tabular}
\end{table}

\begin{figure}[h]
\centering
\includegraphics[width=0.95\textwidth]{07_attribution/shap_if_ocsvm_pak.png}
\caption{Per-sensor attribution precision@$5$ across all six models on HAI test, with $\pm 1\sigma$ error bars across three seeds.  Classical baselines (IF, OCSVM) use SHAP; AE-family models use reconstruction-error attribution from \S\ref{ch:results-attribution}.  Random baseline shown as a dashed line per panel.}
\label{fig:attr-shap}
\end{figure}

Both classical baselines clear the random baseline on every process at p@$5$: IF lifts of $1.39\times$ / $1.35\times$ / $2.18\times$ on P1 / P2 / P3, OCSVM lifts of $1.44\times$ / $1.50\times$ / $2.79\times$.  The interesting result is OCSVM: it ties Dense AE on P1 ($0.69$ vs.\ $0.69$), beats every windowed model on P2 ($0.42$ vs.\ $0.09$ for LSTM-AE / USAD and $0.35$ for TranAD), and sits between LSTM-AE / USAD and Dense AE on P3.  Phase~2 ranks OCSVM near the bottom of the HAI detection table (eTaPR $0.170$, fifth of six), so this is the same detection-vs-localization inversion that Dense AE exhibited in Finding 4.1, now appearing in a kernel method too: the two models with the weakest absolute detection on HAI produce the most actionable per-feature attributions across all three attacked processes.

The mechanism for OCSVM's localization quality is plausible: \code{shap.KernelExplainer} on the OCSVM decision function attributes per-feature contribution by perturbing each feature against the $100$-row training-normal background; features whose perturbation moves the kernel-distance decision the most are exactly those that distinguish the test row from the normal manifold the SVM learned.  That signal is per-row and unaveraged over time, like Dense AE's, so it preserves discrete-flag attacks (P2) that the windowed models smear away.  IF's \code{TreeExplainer} attribution is weaker on P2 / P3 because the model itself isolates rows with fewer effective splits when many features look normal; on P1 (sustained large-magnitude attacks across many sensors) IF localizes well, but on the discrete-flag and rare-process attacks it has less to attribute.

\textbf{Caveat.}  KernelExplainer with a $100$-row background is approximate: the qualitative ranking is robust to background size in our setup, but the exact contribution magnitudes are not, and OCSVM's $100$ attack rows per process per seed is enough for ranking but tighter than the AE-family's full attack set ($n_{\text{P1}} \approx 7000$, $n_{\text{P2}} \approx 1300$, $n_{\text{P3}} \approx 410$ per seed).  As with the AE-family results above, p@k is computed at the process granularity HAI provides; per-sensor evaluation would need labelled per-sensor ground truth which HAI 21.03 does not release.

\subsection{Limitations}

\begin{enumerate}
  \item \textbf{Process-level eval is weaker than per-sensor eval.}  HAI 22.04's README has more detailed per-attack descriptions; a hand-curated per-sensor ground truth on 22.04 would strengthen the claim.
  \item \textbf{Dense AE is non-windowed; LSTM-AE, USAD, TranAD are windowed.}  The localization-vs-detection tradeoff is partly confounded with windowing.  A windowed MLP AE would disentangle; out of scope here.
  \item \textbf{Attention rollout is a single-layer approximation.}  A deeper TranAD could give richer signals; deeper TranAD also takes hours per run on our 4\,GB GPU.
  \item \textbf{KernelExplainer is approximate.}  OCSVM SHAP values are computed with a $100$-row background sample and $128$ Monte-Carlo perturbations per row; tighter settings exceeded our compute budget.  Qualitative rankings should be stable but exact contribution magnitudes are not (\S\ref{sec:finding-shap}).
\end{enumerate}

\chapter{Discussion}
\label{ch:discussion}

\section{What the three findings add up to}

Taken together:
\begin{itemize}
  \item \textbf{Detection (Chapter~\ref{ch:results-detection}).}  SOTA windowed models beat classical ones on same-dataset eTaPR; rankings change with metric choice; PA-F1 inflation is real.
  \item \textbf{Transfer (Chapter~\ref{ch:results-transfer}).}  Source-calibrated thresholds don't transfer; target-calibrated thresholds reveal a genuine classical/SOTA split; within-HAI LOPO shows feature-specific blindness.
  \item \textbf{Attribution (Chapter~\ref{ch:results-attribution}).}  Dense AE beats SOTA on localization; attention rollout adds nothing.
\end{itemize}

The single thesis-level story: \textbf{ICS anomaly detection is not a one-number benchmark.}  Different metrics, calibration regimes, and attribution tasks all reshuffle the ranking.  Any paper that picks one (model, metric, threshold) triple to report has under-constrained its own claim.  Reviewers should ask for the full grid; our contribution is a protocol for producing it.

\section{Likely reviewer objections}

\paragraph{``Your transfer projection is lossy.''}
True, by design (Chapter~\ref{ch:results-transfer}, \S1).  We accepted option~1 (per-variable type tagging) as the simplest defensible scheme.  Option~2 (learned projections) and option~3 (within-testbed-only) remained unused, and we report the ceiling-check against within-HAI LOPO: even in the most favourable within-testbed condition, the ceiling is $0.57$ pointwise F1 --- an order of magnitude better than cross-testbed but still far from same-dataset.  The cross-testbed numbers are not artifacts of a bad projection; they are upper-bounded by the problem itself.

\paragraph{``TranAD reproducibility gap.''}
PA-F1 $0.46$ vs.\ Tuli et al.~\cite{tuli2022tranad}'s $0.97$.  Sources enumerated in \S\ref{sec:tranad-repro}: dataset version, subsample, architecture capacity, threshold protocol.  Our TranAD still ranks \#1 on HAI PA-F1 in our setup --- the qualitative claim survives.  The thesis does not depend on matching published absolutes; it depends on comparing models under a consistent protocol.

\paragraph{``USAD collapsed on Morris --- selection bias?''}
USAD eTaPR $0.009$ on Morris (Chapter~\ref{ch:results-detection}) is not cherry-picked; it is what the model's scores produce.  The collapse is reported as a finding about USAD's assumptions on discrete data, not hidden.

\paragraph{``Only one seed for Phase 4.''}
Acknowledged; listed as a limitation.  Phase 1--2 multi-seed numbers showed eTaPR variance $\leq 0.02$ across seeds for the classical models; attribution seed-variance is similar in the multi-seed runs reported in Chapter~\ref{ch:results-attribution}.  The P1/P2/P3 ordering is robust to that magnitude of noise.

\paragraph{``Process-level attribution is weak.''}
Yes; HAI 21.03 does not publish per-sensor attack targets in machine-readable form.  Chapter~\ref{ch:results-attribution}, \S``Limitations'' lists this as the primary weakness and points at HAI 22.04 as a follow-up.

\section{Practical implications for deployment}

\begin{itemize}
  \item \textbf{Calibrate on target-normal.}  Any deployment should collect a clean-operation period from the target plant before activating the detector and calibrate the threshold there.  Source-trained percentile thresholds produce trivial classifiers (\S\ref{sec:finding-src-cal}).
  \item \textbf{Report three metrics.}  Readers should not trust single-metric rankings.  Our recommendation: pointwise F1 for per-timestep claims, PA-F1 for ``at least one timestep of the event was flagged'' (and flag the inflation per Kim et al.~\cite{kim2022towards}), eTaPR~\cite{hwang2022etapr} for event-level detection with coverage weighting.
  \item \textbf{Consider Dense AE for localization even if a SOTA model is the detector.}  A hybrid pipeline --- SOTA for flagging, Dense AE for attribution --- would outperform either alone on operator-facing triage.
\end{itemize}

\section{What we deliberately did not test}

\begin{itemize}
  \item New detector architectures.  The thesis is empirical, not architectural.
  \item SWaT~\cite{goh2017swat} / WADI~\cite{ahmed2017wadi}.  Adding more single-testbed datasets does not change the cross-testbed contribution.
  \item Real-time streaming performance.  Every number is offline batched.
  \item Federated or differentially-private training.  Out of scope.
\end{itemize}

\section{Future work}

\begin{enumerate}
  \item \textbf{HAI 22.04 per-sensor attribution.}  22.04's README contains per-attack target-sensor descriptions that can be parsed into a ground-truth mapping.  Running the Chapter~\ref{ch:results-attribution} eval against that would upgrade the localization claim from process-level to sensor-level.
  \item \textbf{Multi-seed attribution} across all six models to confirm or refute the LSTM-AE $\equiv$ USAD coincidence (Chapter~\ref{ch:results-attribution}, Finding 4.3).
  \item \textbf{Learned projection for schema alignment} (CLAUDE.md \S5.2 option 2) to test whether the cross-testbed ceiling moves under a richer alignment.  We expect modest improvement given the LOPO ceiling argument above, but the experiment is cheap to run.
  \item \textbf{PA-F1 replacement in the ICS AD literature.}  Our data reproduces the Kim et al.~\cite{kim2022towards} critique on ICS benchmarks; an editorial push to eTaPR as the default would be a small but useful community-level contribution.
\end{enumerate}

